\input{../../../../templates/es/preambulo.tex}
\setbeamertemplate{navigation symbols}{}
\setbeamertemplate{footline}[frame number]
\date{}
\begin{document}
\tituloleccion{04}{Cruza}
\begin{frame}{Pregunta de diseño}
\normalsize
Dados dos padres fijos, ¿qué descendientes puede alcanzar un operador? El caso de Planta Física usa 48 decisiones reales, cuatro edificios por doce meses.
\end{frame}
\begin{frame}{Qué se mide}
\normalsize
El alcance es el soporte de la distribución de hijos. Mido hijos distintos, genes nuevos, clones y fracción legal. Una muestra ilustra; enumerar demuestra el tamaño del soporte.
\end{frame}
\begin{frame}{Varios cortes y uniforme}
\normalsize
\begin{center}\includegraphics[width=.84\textwidth,height=.53\textheight,keepaspectratio]{../../figures/descendencia_03_mas_cortes.png}\end{center}
{\small En esta implementación: un corte 10 hijos, dos cortes 12, tres cortes 8, uniforme 64.}
\end{frame}
\begin{frame}{BLX-\(\alpha\)}
\normalsize
Si \(m=\min(a,b)\), \(M=\max(a,b)\) y \(d=M-m\), entonces \(x'\sim U[m-\alpha d,M+\alpha d]\). Para \(\alpha=.5\), salir del intervalo parental tiene probabilidad \(1/2\).
\end{frame}
\begin{frame}{Acotar no basta}
\normalsize
\begin{center}\includegraphics[width=.84\textwidth,height=.53\textheight,keepaspectratio]{../../figures/descendencia_04_mezcla.png}\end{center}
{\small La mezcla calcula genes nuevos y puede salir del dominio. Acotar sólo repara límites por gen.}
\end{frame}
\begin{frame}{Rutas ilegales}
\normalsize
\begin{center}\includegraphics[width=.84\textwidth,height=.53\textheight,keepaspectratio]{../../figures/descendencia_06_la_division.png}\end{center}
{\small Corte y mezcla no construyen rutas legales en esta prueba. La cruza uniforme sólo recupera un padre legal.}
\end{frame}
\begin{frame}{Dos familias}
\normalsize
\begin{center}\includegraphics[width=.84\textwidth,height=.53\textheight,keepaspectratio]{../../figures/descendencia_07_ordenada.png}\end{center}
{\small OX1 funciona para rutas; una cruza aritmética funciona para valores reales. Elegir la familia es parte del modelo.}
\end{frame}
\begin{frame}{Transferencia al reto}
\normalsize
Para las 48 proporciones del caso, una cruza real propone un hijo. El evaluador compartido comprueba límites por gen, cuota anual, mínimo potable y variación mensual. Acotar cada gen no demuestra factibilidad global.
\end{frame}
\end{document}
